\documentclass[11pt, a4paper]{article}

% --- CODIFICA, LINGUA E TIPOGRAFIA ---
\usepackage[utf8]{inputenc}
\usepackage[T1]{fontenc}
\usepackage[italian]{babel}
\usepackage{lmodern}
\usepackage{microtype} 

% --- GEOMETRIA E LAYOUT ---
\usepackage[top=2.5cm, bottom=2.5cm, left=2.5cm, right=2.5cm, headheight=15pt]{geometry}
\usepackage{fancyhdr}
\usepackage{tabularx}
\usepackage{booktabs}
\usepackage{graphicx}
\usepackage{caption}
\usepackage{xcolor}

% --- COLLEGAMENTI E METADATI PDF ---
\usepackage[
    colorlinks=true,
    linkcolor=black,
    urlcolor=blue!70!black,
    citecolor=black,
    pdftitle={Verbale Interno N. 5 - 13oolean},
    pdfauthor={13oolean}
]{hyperref}

% --- DEFINIZIONE METADATI E COMANDI RIUTILIZZABILI ---
\newcommand{\GroupName}{13oolean}
\newcommand{\GroupEmail}{\href{mailto:gruppo13oolean@gmail.com}{gruppo13oolean@gmail.com}}
\newcommand{\DocTitle}{Verbale Riunione Interna \#6}
\newcommand{\DocDate}{2026-11-09}

% --- CONFIGURAZIONE INTESTAZIONI E PIÈ DI PAGINA ---
\pagestyle{fancy}
\fancyhf{}
\lhead{\textbf{\GroupName}}
\chead{\DocTitle}
\rhead{\DocDate}
\cfoot{\thepage}
\renewcommand{\headrulewidth}{0.4pt}
\renewcommand{\footrulewidth}{0pt}

% ==============================================================================
\begin{document}

% ------------------------------------------------------------------------------
% FRONTESPIZIO
% ------------------------------------------------------------------------------
\begin{titlepage}
    \centering
    
    \includegraphics[width=0.45\textwidth]{../../../shared/logo_unipd.png}\par
    \vspace{0.8cm}
    
    {\scshape\large Università degli Studi di Padova\par}
    {\large Corso di Laurea Triennale in Informatica\par}
    {\large Corso di Ingegneria del Software -- A.A. 2026/2027\par}
    
    \vspace{1cm}
    \includegraphics[width=0.35\textwidth]{../../../shared/logo_gruppo.jpeg}\par
    \vspace{0.8cm}

    {\large Gruppo: \textbf{\GroupName}\par}
    {\large Email: \GroupEmail\par}

    \vfill

    {\Huge\bfseries \DocTitle\par}
    \vspace{0.3cm}
    {\large Data: \DocDate\par}
    
    \vspace{1.5cm}
\end{titlepage}

% ------------------------------------------------------------------------------
% INDICE E INIZIO CONTENUTO
% ------------------------------------------------------------------------------
\pagenumbering{Roman}
\tableofcontents
\clearpage

\pagenumbering{arabic}

\section{Informazioni Generali}
\begin{itemize}
    \item \textbf{Luogo / Piattaforma:} Server Discord
    \item \textbf{Data:} \DocDate
    \item \textbf{Orario:} 15:30 -- 17:12
    \item \textbf{Responsabile:} Michele Tosato
    \item \textbf{Verificatore:} Andrea Guerra
    \item \textbf{Scriba:} Gabriele Maria D'Amico
    \item \textbf{Partecipanti:}
        \begin{itemize}
            \item Andrea Guerra
            \item Francesco Cirigliano
            \item Gabriele Maria D'Amico
            \item Giacomo Lazzarotto
            \item Davide Menegaldo
            \item Michele Tosato
        \end{itemize}
\end{itemize}
\clearpage


\section{Ordine del Giorno}
\begin{enumerate}
    \item Scrittura della mail all'azienda proponente per organizzare un incontro di allineamento.
    \item Analisi dei requisiti e valutazione dello stack tecnologico di progetto.
    \item Configurazione di GitHub Projects, gestione del repository e verifica delle Pull Request in sospeso.
    \item Pianificazione delle prossime attività e calendarizzazione della successiva riunione interna.
\end{enumerate}


\section{Resoconto della Riunione}

\subsection{Analisi dello stack tecnologico e avvio dello studio individuale}
Il gruppo ha affrontato il tema della formazione tecnica sugli strumenti e sui linguaggi richiesti dal capitolato C0, evidenziando l'importanza di allineare le competenze di tutti i membri per garantire una corretta rotazione dei ruoli.
A supporto dello studio individuale è stato condiviso un documento riepilogativo che raccoglie le basi teoriche ed esercizi pratici sullo stack previsto: Java, Spring Boot, Angular, gestione delle API, Git, database relazionali (valutando l'impiego di PostgreSQL, SQLite o MariaDB, su cui il gruppo possiede già conoscenze pregresse) ed altri.
Considerata l'ampiezza degli argomenti e le diverse tempistiche individuali, il gruppo ha stabilito di non fissare una scadenza rigida per il completamento dello studio, ma di gestirlo come un'attività continuativa (\textit{work in progress}), dedicando parte delle riunioni periodiche al confronto sui progressi e alla risoluzione di eventuali dubbi. Inoltre, per garantire uniformità operativa, tutti i componenti provvederanno a configurare il proprio ambiente di sviluppo sulla medesima distribuzione Linux (Ubuntu).

\subsection{Analisi dei requisiti}
In merito alla redazione dell'\textit{Analisi dei Requisiti} è emerso che una stesura eccessivamente dettagliata in questa fase iniziale rischia di introdurre scelte architetturali premature e difficili da verificare.
Il gruppo ha pertanto concordato di procedere con una stesura di base più ad alto livello, con l'obbiettico di identificare gli attori principali, i casi d'uso generali e i potenziali casi limite (quali eventuali errori di lettura dell'OCR o fallimenti nelle richieste), riservando i dettagli implementativi al confronto diretto con l'azienda proponente.

\subsection{Scrittura mail all'azienda per l'organizzazione di un incontro}
Al fine di chiarire i dubbi emersi sull'\textit{Analisi dei Requisiti} e definire con maggiore precisione le tecnologie da adottare, il gruppo ha stabilito come prioritario il contatto via email con l'azienda proponente, con l'obbiettivo di:
\begin{itemize}
    \item Concordare le modalità di comunicazione e la cadenza degli incontri periodici di aggiornamento;
    \item Chiedere indicazioni sugli strumenti di coordinamento interno suggeriti e sull'interazione con il team IA aziendale;
    \item Sottoporre i primi quesiti relativi ai casi limite e ai vincoli tecnologici del progetto.
\end{itemize}

\subsection{Organizzazione su GitHub Projects e gestione del repository}
Il gruppo ha deciso di adottare lo strumento GitHub Projects per semplificare l'organizzazione delle attività, la visualizzazione dello stato di avanzamento e l'assegnazione delle issue.
Sono state inoltre verificate le Pull Request aperte relative all'aggiornamento del sito web (sezione documenti in evidenza), alle correzioni dei verbali pregressi, al \textit{Glossario} e alle \textit{Norme di Progetto} (comprensive delle convenzioni per la scrittura dei commit). Per quanto concerne le label di priorità delle issue, si è deciso di riservare l'assegnazione esplicita principalmente ai task ad alta priorità (\texttt{Priority: High}).

\subsection{Pianificazione della prossima riunione}
In chiusura di seduta, il gruppo ha concordato di aggiornarsi in modalità asincrona nella giornata di venerdì per valutare l'eventuale riscontro da parte dell'azienda proponente e di fissare ufficialmente la prossima riunione interna sincrona per lunedì 2026-11-16 alle ore 15:30.


\clearpage
\section{Decisioni Prese}

\begin{table}[htbp]
\centering
\caption{Riepilogo delle decisioni assunte durante la riunione}
\label{tab:decisioni}
\begin{tabularx}{\textwidth}{l X}
\toprule
\textbf{Codice} & \textbf{Decisione} \\
\midrule
DEC-6.1 & Adozione di GitHub Projects per semplificare l'organizzazione del lavoro e l'assegnazione delle issue. \\
DEC-6.2 & Assegnato il compito di stesura del verbale. \\
DEC-6.3 & Fissata la prossima riunione interna per il 2026-11-16 alle ore 15:30. \\
DEC-6.4 & Avvio dello studio individuale delle tecnologie da utilizzare e allineamento dell'ambiente di sviluppo su Ubuntu. \\
DEC-6.5 & Invio di una mail all'azienda proponente per organizzare un incontro volto a chiarire l'\textit{Analisi dei Requisiti} e il coordinamento periodico. \\
DEC-6.6 & Redazione di una versione base ad alto livello dell'\textit{Analisi dei Requisiti}. \\
\bottomrule
\end{tabularx}
\end{table}

\section{Azioni da Svolgere}

\begin{table}[htbp]
\centering
\caption{Elenco delle azioni e task assegnati}
\label{tab:azioni}
\begin{tabularx}{\textwidth}{l l X l c}
\toprule
\textbf{Codice} & \textbf{Rif.} & \textbf{Descrizione Task} & \textbf{Assegnatario} & \textbf{Scadenza} \\
\midrule
T-6.1 & DEC-6.2 & Stesura del verbale relativo alla riunione tenuta il 19-10-2026 & Gabriele & 2026-11-16 \\
T-6.2 & DEC-6.1 & Configurazione e avvio di GitHub Projects & Francesco Cirigliano  & 2026-11-16 \\
T-6.3 & DEC-6.6 & Stesura base dell'\textit{Analisi dei Requisiti} & Tutti & 2026-11-16 \\
T-6.4 & DEC-6.5 & Contattare l'azienda proponente per fissare un incontro & Michele Tosato & 2026-11-13 \\
T-6.5 & DEC-6.4 & Studio dei linguaggi, dello stack tecnologico e setup dell'ambiente di sviluppo & Tutti & In corso \\
\bottomrule
\end{tabularx}
\end{table}

\end{document}